\section{Survey Methodology}
\label{sec:methodology}

\subsection{Scope and inclusion criteria}
A system is included when it performs or represents at least one of: data distribution or ownership; computation partitioning or placement; communication generation; collective insertion; halo or ghost-region construction; global-to-local address conversion; import/export schedule generation; redistribution or resharding; hierarchical parallelism across nodes or devices; distributed execution-plan generation; or distributed-memory performance modelling used by a compiler. \ours{We classify included work as \emph{central} (automatic distributed-memory parallelization; global-array and data-parallel compilers; Fortran~D, HPF and relatives; automatic message-passing generation; distributed dense and sparse tensor compilers; SPMD partitioners; distributed GPU/accelerator compilers; spatial and wafer-scale compilers) or \emph{bridge} (languages whose compilers implement distributed memory; compiler-integrated runtimes such as inspector--executor systems; PGAS compilers; task/region systems where placement is central; distributed-training systems with substantial plan generation; performance models that drive compiler decisions).} Excluded are distributed databases and operating systems, generic job scheduling, routing protocols, hardware-only papers, application scaling studies, shared-memory-only compilers, single-GPU kernel compilers, and pure runtimes without a compiler component. A small number of architecture papers are retained in the catalog solely to date target classes; they are not classified on the compiler axes.

\subsection{Sources and queries}
Bibliographic metadata were taken, in order of preference, from the DOI record (via Crossref), arXiv, official author or institution copies, official project documentation, and trusted bibliographic indexes. No paywall was bypassed; where full text was not available the catalog marks the entry \texttt{paywalled} or \texttt{abstract-only} and no mechanism is inferred from the title. The initial query set (25 phrases such as ``data distribution compiler'', ``communication set generation'', ``inspector executor distributed memory'', ``SPMD partitioner compiler'', ``wafer scale compiler'') is listed in the repository together with a discovery log recording each query, date, hits, and screening decision. \todo{Report per-database hit counts once the full query sweep over ACM DL and IEEE Xplore is logged.}

\subsection{Snowballing}
From every anchor paper we snowball backward (references), forward (citing works, prioritising those that extend, critique, evaluate, or revive the anchor), laterally (same system, project, or authors), across surveys (works cited by several independent surveys are prioritised), and by \emph{gap}: for every empty cell of the taxonomy coverage matrix we search for that combination specifically.

\subsection{Screening and evidence levels}
Each catalog entry carries one of five statuses: \emph{discovered} (metadata only), \emph{screened} (abstract and introduction checked), \emph{core} (clearly relevant), \emph{read} (enough read to summarise accurately), and \emph{verified} (metadata, categorisation, and summary reviewed against the primary source). Only \emph{read} and \emph{verified} entries support mechanism-level claims in this article. A paper is called a survey only if it is a survey or an invited overview; for example, Lee's ``Techniques for compiling programs on distributed memory multicomputers''~\cite{lee1995techniques} is catalogued as a technique paper.

\subsection{Reading strategy and paper cards}
For most papers we read the abstract, introduction, system overview, related work, and conclusion; we read the full paper when it introduced a major system, defines a taxonomy category, supports a key historical claim, has experiments that we compare, or supplies a mechanism central to the synthesis. Every core paper receives a paper card in the repository listing problem, input abstraction, target, decomposition, mapping, communication, compiler/runtime split, evaluation, limitation, lineage, and the exact sections used.

\subsection{Evidence labelling}
Statements in this article are of three kinds: source-supported facts, which carry a citation; survey taxonomy introduced by us, flagged by ``we categorize'' or ``we group''; and inferences, flagged by ``appears'', ``we observe'', or ``we hypothesise'' and listed with their current evidence in the repository's historical-claims file. Conflicts between sources are recorded rather than resolved silently.

\subsection{Limitations}
Full text of many 1988--2000 papers is paywalled; their summaries rely on open author copies where available. Citation counts are not used as a relevance criterion. Dates of era boundaries are provisional and are justified by the dated events listed in the catalog.
